\section{Conclusion}
PACE combines separately calibrated supervised peer classifiers with a protected local anchor and a heterogeneous bank of anomaly experts for targets that provide benign observations but no attack labels. Across 58 benchmark targets, replacing a seven-detector mean local branch improves target-macro detection in all four primary configurations. The gains range from approximately 2.32 to 2.83 percentage points; at nominal $\alpha=0.05$, they accompany lower mean realized false-alert rates. This is the strongest evidence for the complete local-branch replacement. Rarity-family removal produces the largest detection loss among the tested expert-family ablations, and local-exclusive coverage identifies where local experts add detections beyond the peer union. Separately retained peer decisions also improve detection over mean or maximum peer fusion, with a higher mean false-alert rate.

The contribution is an empirically evaluated architecture for complementary peer and local detection under benign-only target fitting and calibration. Gains are concentrated in TON\_IoT, are smaller in several other strata, and do not consistently increase with larger local fitting budgets. The cohort also informed development, and heterogeneous realized false-alert rates limit operational interpretation. The results support the measured architectural tradeoffs; validation on untouched traffic is needed to establish their generality and acceptable false-alert behavior in deployment.
